% !TeX root = ../main.tex

\section*{Homework \#12 [2025-11-25]}

\begin{problem}[Time-independent version of Wick's theorem]
  Make the canonical transformation to particles and holes for fermions
  $c_{k\lambda} = \theta(k - k_F) a_{k\lambda}
                + \theta(k_F - k) b_{-k\lambda}^\dagger$.
  By applying Wick's theorem, prove the relation
  \begin{multline*}
    c_1^\dagger c_2^\dagger c_4c_3
  = N(c_1^\dagger c_2^\dagger c_4c_3) + \theta(k_F - k_2)
    [\delta_{24}N(c_1^\dagger c_3) - \delta_{23}N(c_1^\dagger c_4)]\\
  + \theta(k_F - k_1) [\delta_{13}N(c_2^\dagger c_4)
  - \delta_{14}N(c_2^\dagger c_3)]
  + \theta(k_F - k_1) \theta(k_F - k_2)
    [\delta_{13}\delta_{24} - \delta_{14}\delta_{23}],
  \end{multline*}
  where the normal-ordered products on the right side now refer to the new
  particle and hole operators, and the subscripts indicate the quantum
  numbers $(k, \lambda)$.
\end{problem}
\begin{solution}
  
\end{solution}

\begin{problem}
  Consider a gas of particles with interaction
  \[
    \hat V = \frac12 \sum_{\bm k k' \bm q \sigma \sigma'} V_q
      c_{\bm k - \bm q\sigma}^\dagger
      c_{\bm k + \bm q\sigma'}^\dagger
      c_{\bm k' \sigma'} c_{\bm k\sigma}.
  \]
  \begin{enumext}
    \item Let $\ket|\phi>$ represent a filled Fermi sea, i.e. the ground state
    of the non-interacting problem. Use Wick's theorem to evaluate an expression
    for the expectation value of the interaction
    energy $\braket<\phi|\hat V|\phi>$ in the non-interacting ground state. Give
    a physical interpretation of the two terms that arise and draw the
    corresponding Feynman diagrams.
    \item Suppose $\ket|\tilde\phi>$ is the full ground state of the interacting
    system. If we add the interaction
    energy $\braket<\tilde\phi|\hat V|\tilde\phi>$ to the non-interacting
    ground-state energy, do we obtain the full ground-state energy? Please
    explain your answer.
    \item (c) Draw the Feynman diagrams corresponding to the second-order
    corrections to the ground-state energy. Without calculation, write out each
    diagram in terms of the electron propagators and interaction $V_q$, being
    careful about minus signs and overall prefactors.
    (If you do not fully understand the question,
    \citep[see][\#7.4.1]{1_coleman2015introduction}.)
  \end{enumext}
\end{problem}
\begin{solution}
  
\end{solution}